\chapter{Mertens's Theorems}

\begin{center}\itshape
Mathematicians have tried in vain to this day to discover some order in the sequence of prime numbers, and we have
reason to believe that it is a mystery into which the human mind will never penetrate.\\
\upshape --- Leonhard Euler (1751)
\end{center}

Picture a machine that draws a whole number at random below a million, and a referee who decides whether it is prime by
trying to divide it by every prime up to a thousand. The prime $2$ catches half the numbers, $3$ a third of them, $p$ a fraction
$\frac1p$. If the referee's tests behaved like independent coin flips, the number would pass them all with probability
$\prod_{p\le1000}\bigl(1-\frac1p\bigr)$, and that should be the proportion of primes. It is a beautiful argument. It is also wrong, by twelve
percent, and it stays wrong by the same twelve percent however far out you go. Finding out exactly how wrong is the business of this chapter,
and the answer involves a constant nobody invited: Euler's $\gamma$, which was born in the harmonic series and has no visible
reason to care about primes.

The story so far. Euler proved that $\sum_p\frac1p$ diverges, and Tome~1 showed that it grows at least like $\log\log x$ (Tome~1, Theorem~2.2).
Chebyshev showed that the number of primes up to $x$ is of the order $\frac x{\log x}$ (Tome~1, Theorem~2.4). Then, in 1874, Franz Mertens, a
professor in Krak\'ow, took nothing more than Chebyshev's bound and the factorisation of $n!$, and found the exact asymptotics of three averages
over primes: $\sum_{p\le x}\frac{\log p}p$, $\sum_{p\le x}\frac1p$, and the referee's product $\prod_{p\le x}\bigl(1-\frac1p\bigr)$. The third is where $\gamma$
sneaks in, and the last section explains how: the product is the reciprocal of the zeta function truncated to the primes up to $x$, and $\gamma$
sits in the expansion of $\zeta$ near $1$.

These theorems are the elementary layer of the theory of primes. They do not give the prime number theorem, but they already
show that if $\pi(x)\frac{\log x}x$ has a limit, that limit is $1$. And they settle the referee's case, which the think box invites you to try first.

\begin{think}
An integer $n$ is prime when no prime $p\le\sqrt n$ divides it. A random integer is divisible by $p$ with probability $\frac1p$. If
these events were independent, the ``probability'' that $n$ is prime would be $\prod_{p\le\sqrt n}\bigl(1-\frac1p\bigr)$. Compute this product for
$n=10^6$ and compare it with the proportion of primes near $10^6$, which is about $\frac1{\log10^6}\approx0.0724$.
\end{think}

\section{Chebyshev's functions}

Counting primes one by one is awkward: every prime counts $1$, whether it is $2$ or a prime with a hundred digits. Chebyshev's trick, in 1852,
was to give each prime a weight $\log p$, roughly its number of digits. With that weighting the prime factors of $n$ add up to $\log n$, and the
primes start to speak the language of factorials, which we understand very well.

Throughout this tome $p$ denotes a prime, $\log$ is the natural logarithm, and $\Lambda$ is von Mangoldt's function:
$\Lambda(n)=\log p$ if $n=p^k$ with $k\ge1$, and $\Lambda(n)=0$ otherwise. Chebyshev's functions are
\[
\begin{gathered}\theta(x)=\sum_{p\le x}\log p,\qq\psi(x)=\sum_{n\le x}\Lambda(n)=\sum_{p^k\le x}\log p.\end{gathered}
\]
Since $\sum_{d\mid n}\Lambda(d)=\log n$ (the prime powers dividing $n$ add up the logarithms of its prime factors with multiplicity),
$\Lambda$ is the natural weight on prime powers: $\Lambda*\mathbf1=\log$ in the notation of Tome~1, Chapter~4.

We use summation by parts in the following form; Tome~4, Proposition~1.7 is the discrete version.

\begin{lemma}[Abel summation]\label{lem:abelint}
Let $(a_n)_{n\ge1}$ be complex numbers, $A(t)=\sum_{n\le t}a_n$, $1\le y<x$, and $f$ of class $C^1$ on $[y,x]$. Then
\[
\begin{aligned}\sum_{y<n\le x}a_nf(n)&=A(x)f(x)-A(y)f(y)\eqbr-\int_y^xA(t)f'(t)\dd t.\end{aligned}
\]
\end{lemma}
\begin{proof}
For $n\le x$ let $n^+=\max(n,y)$. Then $f(x)-f(n^+)=\int_y^x\mathbf1_{n\le t}f'(t)\dd t$, and summing against $a_n$ (a finite sum),
\[
\sum_{n\le x}a_n\bigl(f(x)-f(n^+)\bigr)=\int_y^xA(t)f'(t)\dd t.
\]
The left side is $A(x)f(x)-A(y)f(y)-\sum_{y<n\le x}a_nf(n)$.
\end{proof}

\begin{proposition}[Chebyshev's bounds]\label{prop:chebyshev}
For $x\ge2$:
\begin{enumerate}[label=(\alph*)]
\item $\theta(x)\le x\log4$ and $\psi(x)\le x\log4+2\sqrt x\log x$; in particular $\psi(x)=O(x)$;
\item $\pi(x)\le\sqrt x+\frac{4\log4\cdot x}{\log x}$; in particular $\pi(x)=O\bigl(\frac x{\log x}\bigr)$;
\item $\displaystyle\sum_{n\le x}\log n=x\log x-x+r(x)$ with $|r(x)|\le\log x+1$.
\end{enumerate}
\end{proposition}
\begin{proof}
(a) Tome~1, Theorem~2.4(i) says $\prod_{p\le x}p<4^x$, which is $\theta(x)<x\log4$. Next, $\psi(x)-\theta(x)=\sum_{k\ge2}\theta(x^{1/k})$, and the terms vanish
for $x^{1/k}<2$, that is for $k>\frac{\log x}{\log2}$. Each term is at most $x^{1/2}\log4$, so the sum is at most $\frac{\log x}{\log2}\sqrt x\log4=2\sqrt x\log x$.

(b) The primes in $]\sqrt x,x]$ each contribute at least $\frac12\log x$ to $\theta(x)$, so there are at most $\frac{2\theta(x)}{\log x}$ of them, and at most $\sqrt x$
primes below $\sqrt x$.

(c) Let $N=\lfloor x\rfloor$. Since $\log$ increases, $\int_1^N\log t\dd t\le\sum_{n=2}^N\log n\le\int_1^N\log t\dd t+\log N$, and $\int_1^N\log t\dd t=N\log N-N+1$. The function
$t\mapsto t\log t-t$ has derivative $\log t\le\log x$ on $[N,x]$, an interval of length less than $1$, so $0\le(x\log x-x)-(N\log N-N)\le\log x$. Combine.
\end{proof}

Figure~\ref{fig:chebyshev} shows $\theta(x)/x$ and $\psi(x)/x$: they seem to tend to $1$. That is the prime number theorem, in the form
$\psi(x)\sim x$ (Chapter~5). The difference between $\theta$ and $\psi$ is negligible: by the proof above, $\psi(x)-\theta(x)=O(\sqrt x\log x)$.

\begin{figure}[ht]
\centering
\fig{ch01-chebyshev}
\caption{$\psi(x)/x$ (blue) and $\theta(x)/x$ (orange) for $10\le x\le3000$. The prime powers $p^k$ with $k\ge2$ make up the gap, of order $\frac{\sqrt x}x$.
At $x=3000$ the values are $1.0004$ and $0.977$.}
\label{fig:chebyshev}
\end{figure}

\section{Mertens's first theorem}

The proof of the next theorem is one idea: write $N!$ in two ways. As a product of the numbers up to $N$, its logarithm is about $N\log N-N$. As a
product of primes, its logarithm is a sum over primes, each counted as often as it divides $N!$. Set the two equal and read off the primes.

\begin{theorem}[Mertens, 1874]\label{thm:mertens1}
For $x\ge2$,
\[
\begin{gathered}\sum_{n\le x}\frac{\Lambda(n)}n=\log x+O(1),\qq\sum_{p\le x}\frac{\log p}p=\log x+O(1).\end{gathered}
\]
\end{theorem}
\begin{proof}
Let $N=\lfloor x\rfloor$. By Legendre's formula (Tome~1, Exercise~1.3), the exponent of $p$ in $N!$ is $\sum_{k\ge1}\lfloor N/p^k\rfloor$, so
\[
\log N!=\sum_{p^k\le N}\Bigl\lfloor\frac N{p^k}\Bigr\rfloor\log p=\sum_{n\le x}\Lambda(n)\Bigl\lfloor\frac xn\Bigr\rfloor,
\]
since $\lfloor N/m\rfloor=\lfloor x/m\rfloor$ for every integer $m\ge1$. Write $\lfloor\frac xn\rfloor=\frac xn-\{\frac xn\}$ with $0\le\{\frac xn\}<1$:
\[
\begin{gathered}\log N!=x\sum_{n\le x}\frac{\Lambda(n)}n-\rho(x),\qq0\le\rho(x)\le\psi(x).\end{gathered}
\]
By Proposition~\ref{prop:chebyshev}, $\log N!=x\log x-x+O(\log x)$ and $\psi(x)=O(x)$. Dividing by $x$ gives the first statement.

The two sums differ by the prime powers $p^k\le x$ with $k\ge2$, which contribute at most
$\sum_p\log p\sum_{k\ge2}p^{-k}=\sum_p\frac{\log p}{p(p-1)}\le\sum_{n\ge2}\frac{\log n}{n(n-1)}<\infty$.
\end{proof}

Write $R(x)=\sum_{p\le x}\frac{\log p}p-\log x$, a bounded function. The numerical values $R(10^3)=-1.298$, $R(10^5)=-1.329$, $R(10^7)=-1.3324$
suggest that $R(x)$ converges. It does, to $-\gamma-\sum_p\frac{\log p}{p(p-1)}=-1.33258\ldots$, but this needs the prime number theorem (Exercise~5.2).
Mertens's theorem already says something about the prime number theorem.

\begin{corollary}[Chebyshev, 1852]\label{cor:chebratio}
$\displaystyle\liminf_{x\to\infty}\frac{\theta(x)}x\le1\le\limsup_{x\to\infty}\frac{\theta(x)}x$. So if $\frac{\theta(x)}x$, or equivalently $\pi(x)\frac{\log x}x$, has a limit, the limit is $1$.
\end{corollary}
\begin{proof}
Lemma~\ref{lem:abelint} with $a_p=\log p$ (and $a_n=0$ for $n$ not prime), $A=\theta$, $f(t)=\frac1t$ and $y=1$ gives
\[
\sum_{p\le x}\frac{\log p}p=\frac{\theta(x)}x+\int_1^x\frac{\theta(t)}{t^2}\dd t.
\]
If $\theta(t)\ge(1+\varepsilon)t$ for all $t\ge t_0$, the right side is at least $(1+\varepsilon)\log x-C$, contradicting Theorem~\ref{thm:mertens1}. So the $\limsup$
is at least $1$, and the same argument with $\theta(t)\le(1-\varepsilon)t$ gives the $\liminf$. For the equivalence with $\pi$: $\theta(x)\le\pi(x)\log x$, and
$\theta(x)\ge(\pi(x)-\pi(x^{1-\varepsilon}))(1-\varepsilon)\log x\ge(1-\varepsilon)\pi(x)\log x-x^{1-\varepsilon}\log x$ for every $\varepsilon\in\left]0,1\right[$.
\end{proof}

\section{The sum of the reciprocals of the primes}

Euler knew that $\frac12+\frac13+\frac15+\frac17+\cdots$ diverges. Mertens found out how fast, and the answer is: about as slowly as a series can
diverge and still be called divergent with a straight face.

\begin{theorem}[Mertens, 1874]\label{thm:mertens2}
There is a constant $M$ such that, for $x\ge2$,
\[
\sum_{p\le x}\frac1p=\log\log x+M+O\Bigl(\frac1{\log x}\Bigr),
\]
namely $M=1-\log\log2+\int_2^\infty\frac{R(t)}{t\log^2t}\dd t=0.2614972128\ldots$, the \emph{Meissel--Mertens constant}.
\end{theorem}
\begin{proof}
Let $A(t)=\sum_{p\le t}\frac{\log p}p=\log t+R(t)$. Lemma~\ref{lem:abelint} with $f(t)=\frac1{\log t}$, $f'(t)=-\frac1{t\log^2t}$ and $y=2$ gives
$\sum_{2<p\le x}\frac1p=\frac{A(x)}{\log x}-\frac{A(2)}{\log2}+\int_2^x\frac{A(t)}{t\log^2t}\dd t$. Since $\frac{A(2)}{\log2}=\frac12$ is the term $p=2$,
\[
\sum_{p\le x}\frac1p=1+\frac{R(x)}{\log x}+\int_2^x\frac{\mathrm dt}{t\log t}+\int_2^x\frac{R(t)}{t\log^2t}\dd t.
\]
The first integral is $\log\log x-\log\log2$. Since $|R|\le C$, the second converges as $x\to\infty$, and $\int_x^\infty\frac{|R(t)|}{t\log^2t}\dd t\le\frac C{\log x}$. So the sum is
$\log\log x+M+O\bigl(\frac1{\log x}\bigr)$. The numerical value is computed in the notebook (Exercise~\ref{ex:mertensconst} gives a faster route).
\end{proof}

\begin{figure}[ht]
\centering
\fig{ch01-mertens}
\caption{$\sum_{p\le x}\frac1p-\log\log x$ for $10\le x\le10^7$, in logarithmic scale. It settles at $M=0.26149\ldots$ (dashed); at $x=10^7$ the difference is
$0.261507$, within $10^{-5}$ of $M$.}
\label{fig:mertens}
\end{figure}

So the primes up to $10^{100}$ have reciprocals adding up to about $\log(100\log10)+0.26\approx5.7$: the series diverges, but so slowly that no
computation will ever see it.

\section{Mertens's product and Euler's constant}

Now the referee's product. Taking logarithms turns it into the sum of the previous section plus a convergent correction, so it is of order
$\frac1{\log x}$ up to a constant factor. The surprise is the factor.

\begin{theorem}[Mertens, 1874]\label{thm:mertens3}
As $x\to\infty$,
\[
\prod_{p\le x}\Bigl(1-\frac1p\Bigr)=\frac{e^{-\gamma}}{\log x}\Bigl(1+O\Bigl(\frac1{\log x}\Bigr)\Bigr),
\]
where $\gamma=0.5772\ldots$ is Euler's constant and $e^{-\gamma}=0.5614\ldots$
\end{theorem}

Taking logarithms, $-\log\prod_{p\le x}(1-\frac1p)=\sum_{p\le x}\frac1p+\sum_{p\le x}h(p)$ with $h(p)=-\log\bigl(1-\frac1p\bigr)-\frac1p$. For $0\le u\le\frac12$,
$0\le-\log(1-u)-u=\sum_{m\ge2}\frac{u^m}m\le\frac{u^2}{2(1-u)}\le u^2$, so $H=\sum_ph(p)$ converges, with $\sum_{p>x}h(p)\le\sum_{n>x}\frac1{n^2}\le\frac1{x-1}$. With
Theorem~\ref{thm:mertens2},
\[
\begin{aligned}-\log\prod_{p\le x}\Bigl(1-\frac1p\Bigr)&=\log\log x+M+H\eqbr+O\Bigl(\frac1{\log x}\Bigr).\end{aligned}
\]
So the theorem says exactly that $M+H=\gamma$, and exponentiating gives the stated form, since $e^{O(1/\log x)}=1+O(\frac1{\log x})$. The constants $M$ and $H$
come from elementary sums; $\gamma$ comes from the harmonic series. The link is the zeta function near $s=1$.

\begin{lemma}\label{lem:mplush}
$M+H=\gamma$.
\end{lemma}
\begin{proof}
Let $\sigma=1+\delta$ with $0<\delta\le1$, and compute $\log\zeta(\sigma)$ in two ways as $\delta\to0^+$.

\emph{First way.} Tome~4, Exercise~10.4 gives $\log\zeta(\sigma)=\log\frac1\delta+O(\delta)$.

\emph{Second way.} Taking logarithms in the Euler product (Tome~4, Theorem~10.3),
$\log\zeta(\sigma)=\sum_p-\log(1-p^{-\sigma})=\sum_pp^{-\sigma}+\sum_ph_\sigma(p)$, where $h_\sigma(p)=-\log(1-p^{-\sigma})-p^{-\sigma}$. Since $u\mapsto-\log(1-u)-u$
increases on $[0,1[$, $0\le h_\sigma(p)\le h(p)$, and $h_\sigma(p)\to h(p)$; the series $\sum_ph(p)$ converges, so $\sum_ph_\sigma(p)\to H$ (the series converges uniformly
in $\sigma$ by the Weierstrass test).

It remains to show that $\sum_pp^{-\sigma}=\log\frac1\delta-\gamma+M+o(1)$. Let $S(t)=\sum_{p\le t}\frac1p$, which is $0$ for $t<2$. Lemma~\ref{lem:abelint} with $a_p=\frac1p$ and
$f(t)=t^{-\delta}$ gives $\sum_{p\le x}p^{-\sigma}=S(x)x^{-\delta}+\delta\int_2^xS(t)t^{-1-\delta}\dd t$. As $x\to\infty$, $S(x)x^{-\delta}\to0$ because $S(x)=O(\log\log x)$, so
\[
\sum_pp^{-\sigma}=\delta\int_2^\infty S(t)\,t^{-1-\delta}\dd t.
\]
By Theorem~\ref{thm:mertens2}, $S(t)=\log\log t+M+E(t)$ with $|E(t)|\le\frac C{\log t}$. We treat the three terms, substituting $u=\delta\log t$, so that
$t^{-\delta}=e^{-u}$ and $\frac{\mathrm dt}t=\frac{\mathrm du}\delta$.
\begin{itemize}
\item $\delta\int_2^\infty Mt^{-1-\delta}\dd t=M2^{-\delta}\to M$.
\item The term $J=\delta\int_2^\infty\log\log t\;t^{-1-\delta}\dd t$ equals, since $\log\log t=\log u+\log\frac1\delta$,
\[
J=2^{-\delta}\log\frac1\delta+\int_{\delta\log2}^\infty e^{-u}\log u\dd u.
\]
The first term is $\log\frac1\delta+O(\delta\log\frac1\delta)$, and the integral tends to $\int_0^\infty e^{-u}\log u\dd u$, which is $\Gamma'(1)=-\gamma$ (Tome~4, Exercise~5.7).
\item For the error term,
\[
\begin{aligned}\Bigl|\delta\int_2^\infty E(t)t^{-1-\delta}\dd t\Bigr|&\le C\delta\int_{\delta\log2}^\infty\frac{e^{-u}}u\dd u\\&\le C\delta\Bigl(\log\frac1{\delta\log2}+1\Bigr),\end{aligned}
\]
which tends to $0$.
\end{itemize}
Adding, $\sum_pp^{-\sigma}=\log\frac1\delta-\gamma+M+o(1)$. Comparing the two ways, $\log\frac1\delta+o(1)=\log\frac1\delta-\gamma+M+H+o(1)$, so $M+H=\gamma$.
\end{proof}

The proof of Theorem~\ref{thm:mertens3} is complete. Numerically $H=0.3157184520\ldots$, and $M+H=0.5772156649\ldots$, as it must.

\begin{beauty}
The answer to the think box. For $n=10^6$, $\prod_{p\le1000}(1-\frac1p)\approx\frac{e^{-\gamma}}{\log1000}=0.0813$, while the density of primes near $10^6$ is $0.0724$.
In general the naive product is $\frac{2e^{-\gamma}}{\log n}=\frac{1.1229}{\log n}$, while the prime number theorem says the density is $\frac1{\log n}$: the
``independence'' heuristic overestimates the primes by $12\%$, forever. Divisibility by small primes behaves independently (Exercise~\ref{ex:sieveprob}),
but divisibility by primes near $\sqrt n$ does not, within $[1,n]$: an integer below $n$ has at most one prime factor above $\sqrt n$, and that constraint
is invisible to a product of probabilities. This factor $2e^{-\gamma}$ is the first sign of the difficulty that sieves meet in Chapter~8, and it is why
Cram\'er's random model of the primes (Chapter~10) uses $\frac1{\log n}$ directly instead of a sieve.
\end{beauty}

\section{Exercises}

\begin{exercise}[$\star$]\label{ex:sieveprob}
Fix $y\ge2$, and let $P=\prod_{p\le y}p$. Show that the number of $n\le x$ with no prime factor $\le y$ is $\sum_{d\mid P}\mu(d)\lfloor\frac xd\rfloor$, and that the
proportion of such $n$ tends to $\prod_{p\le y}(1-\frac1p)$ as $x\to\infty$.
\end{exercise}

\begin{exercise}[$\star$]\label{ex:mertensconst}
Show that $M=\gamma+\sum_p\bigl(\log(1-\frac1p)+\frac1p\bigr)$.
\end{exercise}

\begin{exercise}[$\star\star$]\label{ex:chebyshevlower}
Show that $\log\lfloor x\rfloor!-2\log\lfloor x/2\rfloor!\le\psi(x)$, and deduce that $\psi(x)\ge x\log2-3\log x-3$ for $x\ge2$.
\end{exercise}

\begin{exercise}[$\star\star$]\label{ex:averageomega}
Let $\omega(n)$ be the number of distinct prime factors of $n$. Show that
\[
\sum_{n\le x}\omega(n)=x\log\log x+Mx+O\Bigl(\frac x{\log x}\Bigr).
\]
\end{exercise}

\begin{exercise}[$\star\star$]\label{ex:plusproduct}
Show that $\prod_{p\le x}\bigl(1+\frac1p\bigr)\sim\frac{6e^\gamma}{\pi^2}\log x$.
\end{exercise}

\begin{exercise}[$\star\star$]\label{ex:mobiusmean}
Show that $\sum_{n\le x}\mu(n)\lfloor\frac xn\rfloor=1$ for $x\ge1$, and deduce that $\bigl|\sum_{n\le x}\frac{\mu(n)}n\bigr|\le1$.
\end{exercise}

\section{Hints}

\begin{hint}{ex:sieveprob}
Inclusion--exclusion over the primes $p\le y$; the error in each term is less than $1$.
\end{hint}
\begin{hint}{ex:mertensconst}
This is Lemma~\ref{lem:mplush} rewritten.
\end{hint}
\begin{hint}{ex:chebyshevlower}
Compare $\lfloor\frac xn\rfloor-2\lfloor\frac x{2n}\rfloor$ with $0$ and $1$, then use Proposition~\ref{prop:chebyshev}(c).
\end{hint}
\begin{hint}{ex:averageomega}
Swap the sums: $\sum_{n\le x}\omega(n)=\sum_{p\le x}\lfloor\frac xp\rfloor$.
\end{hint}
\begin{hint}{ex:plusproduct}
$1+\frac1p=\frac{1-1/p^2}{1-1/p}$.
\end{hint}
\begin{hint}{ex:mobiusmean}
$\sum_{n\le x}\mu(n)\lfloor\frac xn\rfloor=\sum_{m\le x}\sum_{n\mid m}\mu(n)$.
\end{hint}

\section{Solutions}

\begin{solution}{ex:sieveprob}
\idea Count the complement with inclusion--exclusion: the multiples of $d$ up to $x$ number $\lfloor\frac xd\rfloor$, and $\mu$ supplies the signs. With $y$ fixed there are
only $2^{\pi(y)}$ terms, a constant.
\writeup An integer $n$ has no prime factor $\le y$ exactly when $\gcd(n,P)=1$, and $\sum_{d\mid\gcd(n,P)}\mu(d)=[\gcd(n,P)=1]$ (Tome~1, Chapter~4). Summing over $n\le x$ and
swapping, the count is $\sum_{d\mid P}\mu(d)\#\{n\le x:d\mid n\}=\sum_{d\mid P}\mu(d)\lfloor\frac xd\rfloor$. Replacing $\lfloor\frac xd\rfloor$ by $\frac xd$ costs less than $1$ per term, so the count is
$x\sum_{d\mid P}\frac{\mu(d)}d+O(2^{\pi(y)})=x\prod_{p\le y}(1-\frac1p)+O(2^{\pi(y)})$, expanding the product. Divide by $x$ and let $x\to\infty$. By Theorem~\ref{thm:mertens3},
the limit is about $\frac{e^{-\gamma}}{\log y}$ for large $y$. The error term $2^{\pi(y)}$ is why this argument fails for $y=\sqrt x$, as the beauty box explains.
\end{solution}

\begin{solution}{ex:mertensconst}
\idea Just rearrange the definition of $H$.
\writeup By Lemma~\ref{lem:mplush}, $M=\gamma-H=\gamma-\sum_p\bigl(-\log(1-\frac1p)-\frac1p\bigr)$. Since the series converges fast (its tail after $x$ is below $\frac1{x-1}$),
this computes $M$ much better than Theorem~\ref{thm:mertens2}; the notebook gets ten digits from the primes up to $10^7$, and many more with the identity
$\sum_pp^{-m}=\sum_{k\ge1}\frac{\mu(k)}k\log\zeta(km)$.
\end{solution}

\begin{solution}{ex:chebyshevlower}
\idea The ratio $\frac{(2m)!}{m!^2}$ is the central binomial coefficient that Chebyshev used: it is large, about $4^m$, and all its prime factors are small, so
it forces many primes up to $2m$.
\writeup By the proof of Theorem~\ref{thm:mertens1}, $\log\lfloor x\rfloor!-2\log\lfloor x/2\rfloor!=\sum_{n\le x}\Lambda(n)\bigl(\lfloor\frac xn\rfloor-2\lfloor\frac x{2n}\rfloor\bigr)$. For real $u$,
$\lfloor2u\rfloor-2\lfloor u\rfloor\in\{0,1\}$; with $u=\frac x{2n}$, each bracket is $0$ or $1$, so the sum is at most $\psi(x)$. By Proposition~\ref{prop:chebyshev}(c), the
left side is $(x\log x-x)-2\bigl(\frac x2\log\frac x2-\frac x2\bigr)+r(x)-2r(\frac x2)=x\log2+r(x)-2r(\frac x2)$, and $|r(x)-2r(\frac x2)|\le3\log x+3$ (with the convention
$\log\lfloor x/2\rfloor!=0$ when $x<4$, the bound still holds by a direct check). So $\psi(x)\ge x\log2-3\log x-3$. With Proposition~\ref{prop:chebyshev}(a),
$0.69\le\frac{\psi(x)}x\le1.39$ for large $x$: Chebyshev's theorem, with explicit constants.
\end{solution}

\begin{solution}{ex:averageomega}
\idea Count pairs $(p,n)$ with $p\mid n\le x$ in the other order: each prime $p$ divides $\lfloor x/p\rfloor$ integers up to $x$. Then Mertens's second theorem does
the rest, and Chebyshev bounds the rounding errors.
\writeup $\sum_{n\le x}\omega(n)=\sum_{n\le x}\sum_{p\mid n}1=\sum_{p\le x}\lfloor\frac xp\rfloor=x\sum_{p\le x}\frac1p-\varepsilon(x)$ with $0\le\varepsilon(x)\le\pi(x)=O(\frac x{\log x})$
(Proposition~\ref{prop:chebyshev}(b)). By Theorem~\ref{thm:mertens2}, $x\sum_{p\le x}\frac1p=x\log\log x+Mx+O(\frac x{\log x})$. So a typical integer up to $x$
has about $\log\log x$ prime factors: $3.0$ for $x=10^9$. Chapter~9 shows that most integers have close to that many, and that the fluctuations are Gaussian.
\end{solution}

\begin{solution}{ex:plusproduct}
\idea Turn the product with $+$ signs into Mertens's product, at the cost of a convergent product, which is $\frac1{\zeta(2)}$.
\writeup $\prod_{p\le x}(1+\frac1p)=\prod_{p\le x}(1-\frac1{p^2})\cdot\prod_{p\le x}(1-\frac1p)^{-1}$. The first product tends to $\prod_p(1-p^{-2})=\frac1{\zeta(2)}=\frac6{\pi^2}$ (Euler
product, Tome~4, Theorem~10.3, and $\zeta(2)=\frac{\pi^2}6$). The second is $\sim e^\gamma\log x$ by Theorem~\ref{thm:mertens3}. So the product is
$\sim\frac{6e^\gamma}{\pi^2}\log x=1.0827\ldots\log x$. Expanding, $\prod_{p\le x}(1+\frac1p)=\sum\frac1n$ over the squarefree $n$ with all prime factors $\le x$, which is consistent
with the density $\frac6{\pi^2}$ of the squarefree integers (Chapter~2).
\end{solution}

\begin{solution}{ex:mobiusmean}
\idea $\lfloor\frac xn\rfloor$ counts the multiples of $n$ up to $x$, so the sum counts each $m\le x$ with weight $\sum_{n\mid m}\mu(n)$, which is $1$ for $m=1$ and $0$ otherwise.
The identity then pins $\sum\frac{\mu(n)}n$ down up to the fractional parts.
\writeup $\sum_{n\le x}\mu(n)\lfloor\frac xn\rfloor=\sum_{n\le x}\mu(n)\#\{m\le x:n\mid m\}=\sum_{m\le x}\sum_{n\mid m}\mu(n)=1$, since $\mu*\mathbf1=\varepsilon$ (Tome~1, Theorem~4.4). Write
$\lfloor\frac xn\rfloor=\frac xn-\{\frac xn\}$: then $x\sum_{n\le x}\frac{\mu(n)}n=1+\sum_{n\le x}\mu(n)\{\frac xn\}$. The last sum has $\lfloor x\rfloor$ terms of modulus less than $1$, and the term
$n=1$ is $\{x\}$, so its modulus is at most $\{x\}+\lfloor x\rfloor-1=x-1$. Hence $x\bigl|\sum_{n\le x}\frac{\mu(n)}n\bigr|\le x$. Chapter~5 shows that this sum actually tends to $0$, by the method that proves the prime number theorem, to which this
fact is equivalent.
\end{solution}

\section*{Sources and further reading}

\begin{itemize}
\item F.~Mertens, \emph{Ein Beitrag zur analytischen Zahlentheorie}, Journal f\"ur die reine und angewandte Mathematik 78 (1874), 46--62.
\item T.~M.~Apostol, \emph{Introduction to Analytic Number Theory}, Springer, 1976, Chapters~3 and~4 (averages of arithmetic functions,
Chebyshev's functions, Mertens's first two theorems).
\item G.~Tenenbaum, \emph{Introduction to Analytic and Probabilistic Number Theory}, 3rd ed., AMS Graduate Studies in Mathematics~163, 2015,
Chapter~I.1 (Mertens's theorems, including the product, by a different route).
\item H.~L.~Montgomery and R.~C.~Vaughan, \emph{Multiplicative Number Theory I: Classical Theory}, Cambridge, 2007, Section~2.2 and
Exercise~2.2.1 (the constants $M$ and $\gamma$).
\item A.~Granville, \emph{Harald Cram\'er and the distribution of prime numbers}, Scandinavian Actuarial Journal (1995), 12--28. Discusses the factor
$2e^{-\gamma}$ of the beauty box and why sieve heuristics fail.
\end{itemize}
